% Appendix E --- release-model catalogue: the self-excitation menu of the
% draft's \S6.4 and the partial-release options of its \S6.5, compressed on
% the way in.  These are catalogues; they are presented as catalogues.

\section{Release-model catalogue}
\label{p:app:catalogue}

\Cref{p:sec:odecases} names five ODE cases; \cref{p:sec:selfexcite,p:sec:partial}
state what is settled about self-exciting and partial release. This appendix
writes the cases out and records the candidates behind those statements.

\subsection{Five ODE cases}

Target cells $T$ are fixed; $\Icell$ denotes productive cells, $\Vfree$ free
particles, $E$ eclipse cells, $Q$ the total intracellular stock. Production
into the stock is mean-linear at rate $\nu$ per cell.

\paragraph{Case 1 --- classical BMVR.}
\begin{equation}
\ddt{\Icell}=\gamma T\,\Vfree-d_{\Icell}\Icell,
\qquad
\ddt{\Vfree}=p\,\Icell-c\,\Vfree,
\qquad
R_0=\frac{\gamma T\,p}{c\,d_{\Icell}}.
\label{p:eq:case1}
\end{equation}

\paragraph{Case 2 --- eclipse and classical release.}
\begin{equation}
\ddt{E}=\gamma T\,\Vfree-(\omega+d_E)E,
\qquad
\ddt{\Icell}=\omega E-d_{\Icell}\Icell,
\qquad
\ddt{\Vfree}=p\,\Icell-c\,\Vfree,
\label{p:eq:case2}
\end{equation}
\begin{equation}
R_0=\frac{\gamma T}{c}\cdot\frac{\omega}{\omega+d_E}\cdot\frac{p}{d_{\Icell}}.
\label{p:eq:case2R0}
\end{equation}

\paragraph{Case 3 --- stock and continuous export.}
\begin{equation}
\ddt{\Icell}=\gamma T\,\Vfree-d_{\Icell}\Icell,
\qquad
\ddt{Q}=\nu\Icell-\varepsilon Q-d_{\Icell}Q,
\qquad
\ddt{\Vfree}=\varepsilon Q-c\,\Vfree.
\label{p:eq:case3}
\end{equation}
In the quasi-steady stock regime,
\begin{equation}
Q_*=\frac{\nu}{\varepsilon+d_{\Icell}}\,\Icell,
\qquad
p_{\mathrm{eff}}:=\frac{\varepsilon Q_*}{\Icell}
=\frac{\varepsilon\nu}{\varepsilon+d_{\Icell}},
\label{p:eq:case3qss}
\end{equation}
so classical BMVR is the fast-export limit of stock-and-export dynamics.
Without eclipse,
\begin{equation}
R_0=\frac{\gamma T\,\nu\,\varepsilon}{c\,d_{\Icell}(\varepsilon+d_{\Icell})}.
\label{p:eq:case3R0}
\end{equation}

\paragraph{Case 4 --- stock and proportional full clear.}
\begin{equation}
\ddt{\Icell}=\gamma T\,\Vfree-d_{\Icell}\Icell,
\qquad
\ddt{Q}=\nu\Icell-\delta\frac{Q^2}{\Icell}-d_{\Icell}Q,
\qquad
\ddt{\Vfree}=\delta\frac{Q^2}{\Icell}-c\,\Vfree,
\label{p:eq:case4}
\end{equation}
The mature balance $\nu=\delta\bar x^2+d_{\Icell}\bar x$ gives
\begin{equation}
\bar x=\frac{-d_{\Icell}+\sqrt{d_{\Icell}^2+4\delta\nu}}{2\delta},
\qquad
p_{\mathrm{eff}}=\delta\bar x^2.
\label{p:eq:case4qss}
\end{equation}

\paragraph{Case 5 --- eclipse, stock and proportional full clear.}
\begin{equation}
\ddt{E}=\gamma T\,\Vfree-(\omega+d_E)E,
\qquad
\ddt{\Icell}=\omega E-d_{\Icell}\Icell,
\label{p:eq:case5}
\end{equation}
\begin{equation}
\ddt{Q}=\nu\Icell-\delta\frac{Q^2}{\Icell}-d_{\Icell}Q,
\qquad
\ddt{\Vfree}=\delta\frac{Q^2}{\Icell}-c\,\Vfree,
\qquad
R_0=\frac{\gamma T}{c}\cdot\frac{\omega}{\omega+d_E}\cdot V_\infty.
\label{p:eq:case5R0}
\end{equation}

\subsection{Hataye limiting-dilution numbers}

\begin{table}[htbp]
\centering
\renewcommand{\arraystretch}{1.12}
\small
\begin{tabular}{@{}ll@{}}
\toprule
Quantity & Value \\
\midrule
First detection, bulk production & day $3$; days $3$--$5$ \\
Favoured eclipse stages $n$ & $\approx5$ \\
CD4 division / eclipse death & $\rho_{\rm div}\approx\mu_E\approx0.5/$day \\
Detectable fraction $f_{\rm HIV}$ & $0.41$ \\
Allee threshold $\Lambda$ & $2.3$ latently infected cells \\
RNA copies at that threshold & $5100$ \\
Establishment definition & $\ge2\times10^5$ copies \\
Establishment from one cell & $2\%$ \\
RNA copies at inflection & $\sim5\times10^3$ \\
\bottomrule
\end{tabular}
\caption{Numbers quoted from Hataye et al.\ \cite{hataye2019principles} for
the HIV contrast of \cref{p:sec:hiv}.}
\label{p:tab:hataye}
\end{table}

\subsection{Self-exciting release}

Twelve candidates, four developed below (\cref{p:tab:selfexcite}). They
nest: Model~2, with intensity $\lambda(t)=f(X_t)+\sum_{t_i<t}h(t-t_i)$, is
the general event-history form, and $h\equiv0$ with $f(x)=\varepsilon x$
recovers continuous export (Case~3), $h\equiv0$ with $f(x)=\delta x^2$ and
full clear recovers Model~7 and Cases~4--5, an exponential
$h(t)=\eta\ee^{-\zeta t}$ is equivalent to Model~10, and Model~6 is a coarse
two-level version without a continuous intensity.

\begin{table}[htbp]
\centering
\renewcommand{\arraystretch}{1.1}
\small
\begin{tabular}{@{}clp{0.26\textwidth}cc@{}}
\toprule
\# & Model & Excitation carried by & ODE BMVR? & Renewal? \\
\midrule
1 & Classical Hawkes & event history only & poor & yes (numerical $g$) \\
2 & Hawkes $+$ stock $X$ & events $+$ depletion/refill & via Model 10 & yes \\
3 & Marked Hawkes & size-weighted events & approximate & yes (numerical) \\
4 & ETAS-style & size $+$ time kernel & approximate & yes (numerical) \\
5 & Cox $+$ latent $Z$ & latent excitability & if $Z$ simple & yes \\
6 & ON/OFF telegraph & sustained high phase & yes & yes \\
7 & $\lambda\propto X^2$ & current load only & yes (Cases 4--5) & yes \\
8 & Nucleation biophysics & assembly cascade & hard & yes (numerical) \\
9 & Age since last release & renewal excitation & yes (stages) & yes \\
10 & Jumping rate $\psi(t)$ & shot-noise intensity & yes (mean-field) & yes \\
11 & Spatial Hawkes & local membrane & average first & after coarse-graining \\
12 & Full clear (Cases 4--5) & instant total cascade & yes & yes \\
\bottomrule
\end{tabular}
\caption{Twelve self-exciting release mechanisms. The last two columns are
the practical question: every row admits a renewal description, because
$g(\alpha)=\mean{\lambda(\alpha)\mid\text{alive}}$ always exists; only some
admit a finite ODE one \cite{hawkes1971spectra}.}
\label{p:tab:selfexcite}
\end{table}

\paragraph{Model 7 --- cooperative intensity, $\lambda(x)=\delta x^2$ (or
$\delta x(x-1)$).} No memory of past events, only the current load:
production $\nu$, optional degradation $\mu X$, full-clear dumps at rate
$\delta X^2$. The ODE form is Cases~4--5 under the homogeneous closure, the
renewal form $g(\alpha)=\delta\,\mean{X_\alpha^2\ind\{\text{alive}\}}$.

\paragraph{Model 6 --- ON/OFF telegraph.} The cell, or a membrane site,
switches between OFF and ON at rates $\sigma_{\rm on}$, $\sigma_{\rm off}$,
releasing at rate $\rho$, or $\varepsilon X$, while ON: intermittent surges
without full clear. The simplest ODE version splits the infected class:
\begin{equation}
\ddt{I_{\rm OFF}}=\gamma T\,\Vfree+\sigma_{\rm off}I_{\rm ON}
-(\sigma_{\rm on}+d_{\Icell})I_{\rm OFF},
\label{p:eq:model6}
\end{equation}
\begin{equation}
\ddt{I_{\rm ON}}=\sigma_{\rm on}I_{\rm OFF}
-(\sigma_{\rm off}+d_{\Icell})I_{\rm ON},
\qquad
\ddt{\Vfree}=\rho\,I_{\rm ON}-c\,\Vfree,
\label{p:eq:model6b}
\end{equation}
and quasi-steady switching gives effective budding at the duty cycle,
$p_{\rm eff}=\rho\,\sigma_{\rm on}/(\sigma_{\rm on}+\sigma_{\rm off})$: a
transparent limit back to constant $p$.

\paragraph{Model 10 --- jumping continuous release rate.} A continuous
intensity $\psi(t)\ge0$ jumps by $\eta$ at each release event and decays to
baseline $\psi_0$ at rate $\zeta$,
$\dd\psi=-\zeta(\psi-\psi_0)\dd t+\eta\,\dd N_t$, with flux $\psi(t)$ or
$\psi(t)X_t$: the Markovian, shot-noise form of a Hawkes process with
exponential kernel. Tracking $\Icell$, $Q$ and $\Psi=\sum_i\psi_i$ with
$\bar\psi=\Psi/\Icell$, a mean-field closure gives ODEs for
$(\Icell,Q,\Psi,\Vfree)$ with flux $\approx\bar\psi Q$ --- fading memory
with a finite-dimensional state, at the cost of closures and three extra
parameters.

\paragraph{Model 2 --- Hawkes with stock.} The literal version:
$\lambda(t)=f(X_t)+\sum_{t_i<t}h(t-t_i)$; on an event $X$ drops by the
released amount, between events it refills at rate $\nu$, the cell dies at
rate $d_{\Icell}$. The renewal form
$g(\alpha)=\mean{\lambda(\alpha)\mid\text{alive}}$ is always well defined
and typically numerical; an ODE form exists only after reduction to
exponential $h$, that is to Model~10 (\cref{p:tab:fourmodels}).

\begin{table}[htbp]
\centering
\renewcommand{\arraystretch}{1.1}
\small
\begin{tabular}{@{}>{\raggedright\arraybackslash}p{0.17\textwidth}>{\raggedright\arraybackslash}p{0.19\textwidth}>{\raggedright\arraybackslash}p{0.17\textwidth}>{\raggedright\arraybackslash}p{0.19\textwidth}>{\raggedright\arraybackslash}p{0.17\textwidth}@{}}
\toprule
 & \textbf{7} $\lambda\propto X^2$ & \textbf{6} ON/OFF & \textbf{10} jumping $\psi$ & \textbf{2} Hawkes$+X$ \\
\midrule
Carrier & load only & discrete phase & continuous $\psi$ & event history ($+$load) \\
Memory & Markov in $X$ & Markov in phase & Markov in $\psi$ & non-Markov unless exp.\ $h$ \\
Clusters & high $X\to$ fast dumps & long ON sojourn & high-$\psi$ window & aftershocks \\
ODE BMVR & yes (Cases 4--5) & yes (split $\Icell$) & yes (mean-field $\Psi$) & only via 10 \\
Renewal & yes & yes & yes & yes \\
$\to$ classic $p$ & QSS $\delta\bar x^2$ & duty cycle & QSS $\mean{\psi X}$ & as 10 if exp.\ $h$ \\
\bottomrule
\end{tabular}
\caption{The four developed self-excitation models. Read the last row as
the answer to when each collapses back to a constant release rate.}
\label{p:tab:fourmodels}
\end{table}

\subsection{Partial release}

At load $n$, a release event sends out $B$ with $0\le B\le n$ and retains
$n-B$. The event clock $\lambda(X)$ is free --- constant $\delta$, linear
$\delta X$, or cooperative $\delta X^2$ --- and so is the law of $B\mid X=n$,
for which there are seven options:
\begin{enumerate}
\item[A.] \textbf{Boolean thinning}, the default: each unit released
independently with probability $\varphi$, so
$B\mid X=n\sim\mathrm{Binomial}(n,\varphi)$; $\varphi\to0$ gives tiny
fractional clear, $\varphi=1$ full clear, and generating functions stay
tractable.
\item[B.] Fixed fraction $B=\mathrm{round}(fn)$: A's mean with $\varphi=f$,
less stochasticity. \ \textbf{C.} Random fraction ($U\sim\mathrm{Beta}$, or
uniform thinning): more variable bursts than B.
\item[D.] Fixed capped batch $B=\min(k,n)$: $k=1$ is unit budding in the
frequent-event limit, large $k$ approaches full clear. \ \textbf{E.}
Size-biased release, $\pr{B=j\mid n}\propto j$: natural if one site releases
a random assembled blob.
\item[F.] Threshold-then-partial: events only for $X\ge x_*$, then A or B.
\ \textbf{G.} Short high-export window: a triggered burst of high
$\varepsilon$, partial clear in expectation, related to Model~6.
\end{enumerate}

The single-cell boolean process produces at rate $\nu$, fires release events
at rate $\lambda(X)$ with $B\sim\mathrm{Binomial}(X,\varphi)$ released and
$X\to X-B$, and dies at rate $d_{\Icell}$; its nested limits are
\cref{p:tab:boolean}. Option~A with a load-dependent clock is the
intermediate used in \cref{p:sec:partial}: one transparent parameter between
mild partial clear and full dumping, mechanistic, nested in Cases~3--4 at
the mean-field level, and straightforward to simulate.

\begin{table}[htbp]
\centering
\renewcommand{\arraystretch}{1.1}
\small
\begin{tabular}{@{}ll@{}}
\toprule
Choice & Recovers \\
\midrule
$\varphi=1$, $\lambda=\delta X^2$ & full clear, cooperative (Case~4 style) \\
$\varphi=1$, $\lambda=\delta X$ & full clear, linear hazard \\
$0<\varphi<1$, $\lambda=\delta X$ & the main partial-clear intermediate \\
$B\equiv1$, $\lambda=\varepsilon X$ & unit budding (Case~3, discrete) \\
\bottomrule
\end{tabular}
\caption{What the boolean process recovers at its corners. The middle row
is the only one that is not already somewhere else in the chapter.}
\label{p:tab:boolean}
\end{table}

\FloatBarrier
